\subsection*{Strong cancellation persists in the exact source complement}

The proved ordinary source residual does not mean that significant
prime/non-prime cancellation is confined to one direction. The following
fixed-window certificate tests that possible shortcut on the literal
source, rather than on sampled basis or random vectors.

\begin{proposition}[A certified cancelling direction off the source]
\label{prop:v121-complement-cancellation}
At $\lambda=3$ there exists an inversion-even unit Fourier polynomial
$v\in E_{64}$ exactly orthogonal to the repaired source $p_3$ such that
\begin{equation}\label{eq:v121-complement-cancellation}
 \begin{gathered}
 \ip{v}{p_3}=0,\qquad
 0<QW_3(v,v)<3\cdot10^{-31},\\
 q_{\rm np}(v)>0.4,\qquad q_{\rm pr}(v)<-0.4.
 \end{gathered}
\end{equation}
Here $q_{\rm np}=q_{\rm ar}+q_{\rm po}$ includes both the pure
archimedean term and the pole term, and $q_{\rm pr}$ is the signed
prime contribution, so $QW_3=q_{\rm np}+q_{\rm pr}$.
\end{proposition}
\begin{proof}
Let $V$ have the two real dyadic columns supplied in
Appendix~\ref{app:v121-audit}, in the normalized even basis of $E_{64}$.
Let $c$ be the exact rational source candidate from
Proposition~\ref{prop:v113-source-certificate}, converted to the same
parity coordinates. Set $a=V^*c$ and $x=(-a_2,a_1)^t$. The exact
algebra gives $\ip{Vx}{c}=0$. The certified Gram matrix $V^*V$ is
positive definite and $\norm{Vx}>0$, so put $w=Vx/\norm{Vx}$.
All vectors are real here; the stated Hermitian convention is unchanged.

Outward interval evaluation of the separate finite matrices gives
\[
 \begin{aligned}
 2.5844134\cdot10^{-31}&<QW_3(w,w)<2.5844135\cdot10^{-31},\\
 q_{\rm np}(w)&>0.4173237590,\\
 q_{\rm pr}(w)&<-0.4173237590.
 \end{aligned}
\]
These are inequalities for exact dyadic/rational constructions;
approximate eigenvectors were used only to propose the dyadic columns.

Write $u=P_{64}p_3/\norm{P_{64}p_3}$ and $\widehat c=c/\norm c$.
The existing exact-source certificate, matched to $c$ by its file hash,
gives $\norm{P_u-P_{\widehat c}}<\delta=4\cdot10^{-36}$.
Since $w\perp c$, define
\[
 v=\frac{(I-P_u)w}{\norm{(I-P_u)w}}.
\]
Then $\norm{v-w}<2\delta$. Also $v\in E_{64}$, hence $v\perp u$
implies $v\perp p_3$, not just orthogonality to an approximate source.
For any finite Hermitian matrix $H$,
\[
 |\ip{Hv}{v}-\ip{Hw}{w}|
 \le2\norm H\norm{v-w}<4\delta\norm H.
\]
The certified all-row bounds for the compressed matrices are
\[
 \norm{W_{64}}<8.926,\qquad
 \norm{H_{\rm np,64}}<4.176,\qquad
 \norm{H_{\rm pr,64}}<5.490.
\]
Applying these bounds to the full interval enclosures, before rounding,
gives
\[
 2.5829853\cdot10^{-31}<QW_3(v,v)<2.5858415\cdot10^{-31}
\]
and preserves both strict $0.4$ bounds. A finite-support polynomial's
complete quadratic form equals its finite matrix quadratic form.
No invariant-subspace or full operator-residual assertion is required.
\end{proof}

Thus $QW_3(v,v)/q_{\rm np}(v)<7.5\cdot10^{-31}$ inside the exact
source complement. This rules out an explanation in which substantial
archimedean-plus-pole dominance eliminates the need for cancellation
everywhere off that source. It does not rule out positivity on the
complement, and exhibits neither a negative Weil direction nor an
additional exact null vector. Positive basis-vector or typical
random-vector energies cannot control these coherent linear combinations.
The pole term must be included explicitly: calling $q_{\rm np}$ the
pure archimedean contribution would change the asserted comparison.

